\section{诉的基本理论}
\begin{enumerate}
  \item 理论界内容非常丰富，法考呢考点非常的明确。
  \item \bs{主体当事，客体标的。分类请求，排除非诉。反诉独立。主客合并，单纯、预备和重叠}。
  \item 要素。
        \begin{enumerate}
          \item 主体叫\zd{当事人}。
          \item 客体又叫\zx{诉讼标的}。
                \begin{enumerate}
                  \item 双方当事人发生纠纷并请求法院作出裁判的那种\zx{实体法律关系}。
                  \item 考试的时候就考一个问题，诉讼标的的区分和判断，\bs{不带“关系”的绝对不能选}。
                  \item vs \zx{诉讼请求}。\al{我要解除合同、我要租金 5000}。诉讼标的 \yjt \al{租赁关系}。
                \end{enumerate}
          \item 诉的理由。
        \end{enumerate}
  \item 分类。根据\zx{原告的诉讼请求}。每个诉对应一种判决。\al{确认之诉 \yjt 确认判决}。
        \begin{enumerate}
          \item \zx{确认}之诉。确认某种法律关系\zx{是否存在、有效}。积极的、消极的。。
          \item \zx{给付}之诉。原告要求判令被告\zx{履行义务}（财务、行为 \yjt 作为、不作为）。
          \item \zx{形成/变更}之诉。要求对一个\zx{既存的}法律关系（主体、客体、内容变一个）予以\zx{变更、消灭}。
                实践当中遇到的绝大多数民事纠纷都是形成之诉。\al{离婚、判小孩归自己}。
                \begin{enumerate}
                  \item \bs{解除合同是确认之诉}。胜诉 \yjt \zd{形成力} \yjt 判决生效之日即为法律关系消灭、变更之时。
                        “既存”要求形成的时候这个法律关系\zd{仍存在}，不存在难道要变更空气么？
                  \item \bs{撤销权是典型的形成之诉}。形成权 \yjt 通知即可行使。形成诉权 \yjt 必须诉讼方式。
                \end{enumerate}
                \begin{kk}{小 Tips}
                  \begin{enumerate}
                    \item 诉的分类的\zx{前提是诉及诉讼}案件，\bs{非讼案件不讨论诉的分类}。考试经常用非讼程序干扰你。
                    \item \zd{分类唯一依据是原告请求}，与被告的\zd{答辩意见无关}，是干扰信息。
                    \item 考试最简单的是变客体，然后主体、内容变来变去。
                  \end{enumerate}
                \end{kk}
        \end{enumerate}
  \item \zx{反诉}。在诉讼中被告向原告提出一个\bs{独立}的请求。
        \begin{enumerate}
          \item 构成要件。
                \begin{enumerate}
                  \item 主体。具有\zx{同一性}。一定是本诉的被告，告本诉的原告。
                  \item 诉讼地位。\zx{完全相反}。\bs{选项看见相同就是错的}。
                  \item 本、反诉之间\zx{有牵连关系}。\bs{但不要求基于相同的法律关系}。
                  \item 要在本诉\zx{一审的法庭辩论终结前}（当事人表达意见的最后机会）提出。
                \end{enumerate}
          \item 基于牵连关系，直接向\zd{受理本诉的法院}提（\zd{专属管辖除外}）
                \yjt 被驳回 \ydy 驳回起诉 \yjt 上诉 \yjt 成立 \yjt \zx{撤销原驳回裁定，指定原审法院继续审理}。
          \item 本、反诉必须适用\zx{同一程序}审理。受理费减半。
          \item 二审中提反诉。\zx{调解不成，告知另诉}。见二审程序部分。
        \end{enumerate}
        \begin{kk}{反诉是一个独立的诉}
          \begin{enumerate}
            \item 核心。\zx{不依赖本诉的存在而存在}。找出被告主张 \yjt 去掉原告起诉 \yjt 能否\bs{单独直接}向法院起诉。
            \item 不会因为本诉的消灭而消灭。各判各，判决结果互不影响。
            \item 考你反诉的判断。考你反诉与\zd{反驳}的区分。\zx{反驳 \bdy 抗辩（权） \bdy 反诉}。
                  \begin{enumerate}
                    \item 都可以用来反对对方请求。
                    \item \zx{反驳、反诉互斥}，但抗辩可以变现为反诉，反诉也可以变现为抗辩。
                          \al{要求付租金，抗辩合同已解除 \yjt 可以成立、反诉合同已解除 \yjt 可以成立。\xhx{但绝不可能是反驳}}。
                  \end{enumerate}
            \item \zd{绝对不允许去看法律效果，不能去看“抵销”、“折抵”}。
          \end{enumerate}
        \end{kk}
  \item 分离。合起来太复杂了，就拆开。没啥考点，
  \item 合并。审理可以合并，如果诉讼标的不同则判决必须分开。
        \begin{enumerate}
          \item 主体的合并。\zx{针对人}。一方 \hz 双方存在 2 人以上。\al{共同诉讼}。
          \item 客体的合并。只有 \zd{2 个人}。针对同一被告、同一原告之间的\zd{诉讼标的 \hz 诉讼请求}。
                \begin{enumerate}
                  \item 单纯合并。\zd{诉讼标的之间没有必然的联系}。分别审理，分别判决。\al{借款合同、买卖合同}。
                  \item 预备合并。\zd{对象是诉讼请求}。\bs{顺序}处理。
                        \begin{enumerate}
                          \item \zx{主诉 + 备诉}。\al{平行志愿}。
                          \item 前诉是后诉的\zx{先决}（命题人观点）。\al{离婚诉讼 \yjt 财产分割、子女抚养}。
                        \end{enumerate}
                  \item 重叠/选择合并。诉讼请求相同，\zd{诉讼标的不同}。\zd{择一}。\bs{记这个例子就行}。% https://q.uiver.app/?q=WzAsMixbMCwwLCLlvKDkuIkiXSxbMSwwLCLmnY7lm5siXSxbMCwxLCJcXHlpIOenn+i1geWQiOWQjFxcXFxcXGVyIOaJgOacieadg+WFs+ezuyIsMSx7ImxldmVsIjoyfV0sWzAsMSwi6L+U6L+Y5oi/5bGLIiwxLHsiY3VydmUiOi0xfV1d
\begin{tikzcd}
  \txt{甲} & & & \txt{乙}
  \arrow["\txt{\yi 租赁合同}", shift left=-1, from=1-1, to=1-4]
  \arrow["\txt{返还房屋}", curve={height=-18pt}, from=1-1, to=1-4]
  \arrow["\txt{\er 所有权关系}"', shift right=3, from=1-1, to=1-4]
\end{tikzcd}

                \end{enumerate}
        \end{enumerate}
  \item 诉讼地位。当事人及其他诉讼参与人在诉讼中的法律地位。
  \item 诉讼行为。诉讼中所进行的全部行为活动。
\end{enumerate}
